\documentclass[a4paper]{article}
% Español (México) con XeLaTeX: fontspec para fuentes Unicode, polyglossia
% para la división silábica y los nombres del documento en la variante mexicana.
\usepackage{fontspec}
\usepackage{polyglossia}
\setdefaultlanguage[variant=mexican]{spanish}
% Los nombres chinos van como «pinyin (original)»: xeCJK da a los caracteres
% Han su propia fuente; sin ella XeLaTeX los omite con un aviso.
% FandolSong no cubre algunos caracteres de nombres (U+73FA); WenQuanYi Zen Hei sí.
\usepackage[AutoFallBack=true]{xeCJK}
\setCJKmainfont{FandolSong-Regular.otf}
\setCJKfallbackfamilyfont{\CJKrmdefault}{WenQuanYi Zen Hei}
% polyglossia no traduce los nombres de listings.
\AtBeginDocument{\providecommand{\lstlistingname}{}\renewcommand{\lstlistingname}{Listado}%
  \providecommand{\lstlistlistingname}{}\renewcommand{\lstlistlistingname}{Índice de listados}}
\usepackage{amsmath, amssymb}
\usepackage{graphicx}
\usepackage[margin=2.35cm]{geometry}
\usepackage[most]{tcolorbox}
\usepackage{booktabs}
\usepackage{tabularx}
\usepackage{longtable}
\usepackage{array}
\usepackage{float}
\usepackage{xcolor}
\usepackage{hyperref}
\usepackage{fancyhdr}
\setCJKmainfont[AutoFakeBold=2.4]{AR PL UMing CN}
\setCJKsansfont[AutoFakeBold=2.4]{Droid Sans Fallback}
\setCJKmonofont[AutoFakeBold=2.4]{Droid Sans Fallback}
\hypersetup{colorlinks=true, linkcolor=blue!55!black, urlcolor=blue!60!black}
\graphicspath{{./}{figures/}}
\setlength{\parskip}{0.55em}
\setlength{\parindent}{2em}
\setlength{\emergencystretch}{3em}
\renewcommand{\arraystretch}{1.18}
\newtcolorbox{knowledgebox}[1]{enhanced, colback=blue!5!white, colframe=blue!70!black, colbacktitle=blue!70!black, coltitle=white, fonttitle=\bfseries, title={#1}, sharp corners, breakable}
\newtcolorbox{importantbox}[1]{enhanced, colback=yellow!10!white, colframe=yellow!75!black, colbacktitle=yellow!75!black, coltitle=black, fonttitle=\bfseries, title={#1}, sharp corners, breakable}
\newtcolorbox{warningbox}[1]{enhanced, colback=red!5!white, colframe=red!70!black, colbacktitle=red!70!black, coltitle=white, fonttitle=\bfseries, title={#1}, sharp corners, breakable}
\pagestyle{fancy}
\fancyhf{}
\fancyfoot[C]{\thepage}
\fancyfoot[R]{\footnotesize\color{gray}\href{https://github.com/hqhq1025/ai-course-notes}{AI Course Notes}}
\renewcommand{\headrulewidth}{0pt}
\renewcommand{\footrulewidth}{0pt}
\newcommand{\videourl}{https://www.youtube.com/watch?v=RxXVq7-sJzM}
\newcommand{\repourl}{https://github.com/hqhq1025/ai-course-notes}

\begin{document}
\begin{titlepage}
\centering
{\Large Notas didácticas de una entrevista a fondo\par}
\vspace{1.0cm}
{\huge\bfseries Modelo grande de CEO, VLA, Agent OS y sabiduría humana}\par
\vspace{0.5cm}
{\Large Zhang Xiaojun conversa con Li Xiang (李想), fundador de Li Auto (理想汽车)\par}
\vspace{0.35cm}
{\large Elaboradas a partir del video público de Zhang Xiaojun Podcast, la transcripción hecha con GPU, la estructura de capítulos y diagramas conceptuales propios\par}
\vspace{0.3cm}
{\large Publicado el 2025-10-30; grabado en 2025-04\par}
\vspace{0.85cm}
\includegraphics[width=0.88\textwidth,height=0.42\textheight,keepaspectratio]{cover.jpg}\par
\vfill
\begin{tcolorbox}[width=0.92\textwidth, colback=black!2!white, colframe=black!55, sharp corners]
\textbf{Título del video}: 118. Segunda entrevista de 3 horas con Li Xiang: modelo grande de CEO, MoE, Liang Wenfeng (梁文锋), VLA, energía, memoria, la lucha contra la naturaleza humana, relaciones íntimas, la sabiduría humana\par
\textbf{Canal del video}: Zhang Xiaojun Podcast\par
\textbf{Duración del video}: 02:46:22\par
\textbf{Enlace del video}: \href{\videourl}{\nolinkurl{\videourl}}\par
\textbf{Nota sobre el material}: YouTube no tiene subtítulos disponibles; el texto se elaboró a partir de la transcripción con faster-whisper large-v3 en CUDA, los capítulos del video, la portada y figuras didácticas propias. Las tomas fijas de la entrevista no se repiten en el cuerpo del texto.\par
\textbf{Página del proyecto}: \href{\repourl}{\nolinkurl{\repourl}}\par
\end{tcolorbox}
\end{titlepage}

\tableofcontents
\newpage
\section{Guía de lectura: entrevistar a un CEO como si fuera un gran modelo MoE}

Esta sección establece primero el método de lectura de toda la entrevista. En la apertura, Zhang Xiaojun (张小珺) propone un marco muy útil: tratar a Li Xiang (李想) como un «gran modelo CEO» y suponer que tiene una arquitectura MoE (Mixture of Experts, mezcla de expertos). Las tres primeras rondas invocan, en orden, al experto técnico, al experto en estrategia y al experto en organización; en la segunda mitad, las preguntas avanzan del modelo, el automóvil y los robots hacia la energía, las relaciones íntimas, la memoria y la sabiduría. Este marco no es una broma: convierte una entrevista de tres horas sobre una persona en una reflexión sistemática sobre las empresas, los productos y las personas en la era de la IA.

El objetivo de esta nota no es repetir la transcripción literal, sino organizar la entrevista en una estructura que se pueda estudiar. La primera línea es la técnica de IA: la ventana de contexto humana, las mejores prácticas de DeepSeek, VLA, el modelo del mundo, Agent OS, Action y la alineación. La segunda línea es la estrategia de la empresa: escala, necesidades de los usuarios, productos tecnológicos, capacidades organizacionales y por qué Li Auto (理想) incluye el automóvil, los robots y el sistema operativo en su visión de los terminales de la era de la AGI. La tercera línea trata de las personas y la organización: la energía, las relaciones íntimas, el cerebro común, el corazón común y la definición de «sabiduría» de Li Xiang.

\begin{figure}[H]
\centering
\includegraphics[width=0.96\textwidth]{ceo-moe-interview-map.png}
\caption{Mapa de la entrevista al gran modelo CEO: Li Xiang tratado como una arquitectura MoE, en la que se invocan conjuntamente la técnica, la estrategia, la organización y las personas. Diagrama conceptual propio, elaborado a partir de la conversación en 00:01:56--00:02:20 y 00:36:18--00:36:25.}
\end{figure}

\begin{knowledgebox}{Lectura de la figura: por qué MoE organiza todo el episodio}
La intuición de MoE es que «distintos expertos se encargan de distintas capacidades». La entrevista usa esta metáfora para dividir las respuestas de Li Xiang en varios canales de expertos: el experto técnico explica VLA y el modelo del mundo, el experto en estrategia explica la escala y la ruta de plataforma, el experto en organización explica la densidad de talento y la metodología, y al final se vuelve a la energía, las relaciones y la sabiduría de una persona. Al leer este episodio, lo importante no es memorizar todas las opiniones, sino observar cómo estos canales de expertos se invocan entre sí.
\end{knowledgebox}

\begin{importantbox}{Tesis central de este episodio}
En la era de la IA, una empresa no se limita a conectar un modelo a un producto: tiene que reescribir al mismo tiempo la forma del producto, la estructura organizacional, el ciclo cerrado de acción y las relaciones humanas. El automóvil, los robots, Agent OS, el modelo del mundo y las relaciones íntimas parecen temas dispersos, pero en realidad responden a la misma pregunta: cómo entra la inteligencia en el mundo real y genera valor de forma continua.
\end{importantbox}

\begin{knowledgebox}{Términos clave: conceptos centrales de este episodio}
\begin{tabularx}{\textwidth}{p{0.20\textwidth}p{0.30\textwidth}X}
\toprule
Término & Problema que resuelve & Significado en este episodio \\
\midrule
MoE & Repartir las capacidades del modelo entre varios expertos & Designa la arquitectura de modelos como DeepSeek y también sirve de metáfora de las capacidades de múltiples expertos de un CEO. \\
VLA & Pasar de la visión y el lenguaje a la acción & Vision-Language-Action, la ruta técnica central con la que Li Auto aborda la conducción autónoma y los robots. \\
World Model & Predecir los cambios del entorno & En escenarios de tráfico se usa para evaluar, generar datos de entrenamiento y podría convertirse en el sistema de operación de la conducción autónoma. \\
Agent OS & Expresión alternativa al Agent general & Li Xiang considera que a corto plazo no habrá un único Agent general, pero sí un sistema operativo en el que funcionen Agents especializados. \\
Action & Que la IA cambie realmente el mundo & Cuando la IA invoca herramientas, activos, vehículos y computadoras, la alineación, los permisos y la responsabilidad adquieren más peso. \\
\bottomrule
\end{tabularx}
\end{knowledgebox}

\subsection{Resumen de la sección}

La dificultad de EP118 está en su enorme amplitud. Al desglosarlo con la metáfora MoE, la línea principal se vuelve mucho más clara: el experto técnico responde «cómo entra el modelo en el mundo físico», el experto en estrategia responde «cómo se convierte la empresa en la plataforma de la siguiente generación», y la organización y las personas responden «quién soporta este cambio».
\section{La ventana de contexto humana: el ser humano reduce la entropía, la IA la aumenta}

Esta sección parte de la «conversación de texto largo» con la que abre el episodio. Zhang Xiaojun (张小珺) le pregunta a Li Xiang (李想): si fueras un modelo grande, ¿qué tan grande sería tu ventana de contexto? La respuesta de Li Xiang no sigue la vía de las cifras de tokens, sino que se desplaza hacia la diferencia entre el ser humano y la IA: las personas no son buenas para procesar información muy compleja, así que deben reducir la entropía; la IA es buena para procesar volúmenes enormes de información, lo que se parece más a aumentar la entropía. Aquí, «reducir la entropía» no es la definición estricta de la física, sino un método de gestión y de cognición: comprimir un mundo complejo en unos cuantos juicios clave, unas cuantas acciones y unos cuantos principios.

\begin{figure}[H]
\centering
\includegraphics[width=0.94\textwidth]{human-ai-context-window.png}
\caption{Contexto humano frente a contexto de la IA: al ser humano le conviene más restar; a la IA, procesar información a gran escala. Diagrama conceptual propio, elaborado a partir de la conversación de 00:02:35--00:04:10.}
\end{figure}

\begin{knowledgebox}{Lectura de la figura: una ventana de contexto más grande no hace a nadie más parecido a un CEO}
Un modelo grande puede procesar muchos tokens, pero el valor de un CEO suele estar en la compresión: convertir información compleja en estrategia, organización y acción. A la inversa, si la IA se queda en «saber mucho» y no se conecta con herramientas, acciones y retroalimentación, tampoco puede generar valor real.
\end{knowledgebox}

\subsection{Tres niveles de herramientas: herramientas de información, herramientas de apoyo y herramientas de producción}

La sección anterior distinguió la capacidad de contexto del ser humano y la de la IA; esta sección da un paso más y pregunta cómo se convierte esa diferencia de capacidades en valor de producto. Li Xiang divide las herramientas en tres niveles. Las herramientas de información permiten a las personas encontrar información; las herramientas de apoyo aumentan su eficiencia; las herramientas de producción generan directamente una productividad por la que se paga. Esta clasificación es importante porque ofrece un criterio más firme para valorar los productos de IA: no se trata de si pueden responder, sino de si el usuario está dispuesto a pagar por ellos y a confiarles trabajo real.

\begin{importantbox}{Criterio para identificar una herramienta de producción}
Si una herramienta de IA solo hace que las personas «se sientan inteligentes», puede seguir siendo apenas una herramienta de información; solo cuando puede sustituir o liberar el trabajo real y frecuente de las personas, y el usuario está dispuesto a pagar por el resultado, empieza a entrar en el nivel de herramienta de producción.
\end{importantbox}

\subsection{DeepSeek: aplicar las mejores prácticas de la forma más sencilla}

A continuación, la entrevista se dirige a DeepSeek. Lo que Li Xiang aprendió de DeepSeek no es una técnica puntual, sino «aplicar de la forma más sencilla las mejores prácticas de la humanidad»: primero investigar, luego desarrollar y después convertirlo en negocio; si la brecha es de resultados, se ajusta la estrategia; si la brecha es de capacidades, se vuelve a la construcción de capacidades. Esta ruta parece sencilla, pero va contra la naturaleza humana, porque exige que la organización reconozca una y otra vez que sus capacidades son insuficientes, en lugar de avanzar a la fuerza solo con acciones de negocio.

\begin{figure}[H]
\centering
\includegraphics[width=0.96\textwidth]{deepseek-best-practice.png}
\caption{Las mejores prácticas aprendidas de DeepSeek: usar de la forma más sencilla las mejores prácticas de la humanidad y luego consolidar la organización y la ingeniería. Diagrama conceptual propio, elaborado a partir de la conversación de 00:15:01--00:19:10.}
\end{figure}

\begin{knowledgebox}{Lectura de la figura: DeepSeek como caso de organización y no solo como caso de modelo}
La secuencia de la figura, que va de MoE a la práctica, expresa una ruta de aprendizaje organizacional: primero investigar el problema, luego construir capacidades, después validarlas con ingeniería y, por último, convertir esas capacidades en resultados de negocio. A Li Xiang le interesa DeepSeek porque muestra cómo un equipo pequeño, una alta densidad de talento, el código abierto y la práctica de ingeniería pueden sumar fuerzas.
\end{knowledgebox}

\begin{warningbox}{Las «mejores prácticas» no producen automáticamente los mejores resultados}
Las mejores prácticas suelen ir contra la naturaleza humana, porque exigen bajar el ritmo para investigar, construir capacidades y hacer revisiones retrospectivas. Actuar según el impulso satisface más a la naturaleza humana, pero facilita confundir el problema con uno de estrategia de corto plazo y pasar por alto la verdadera brecha de capacidades.
\end{warningbox}

\subsection{Por qué desarrollar de todos modos capacidades propias de modelo base}

En la entrevista hay un detalle: dentro de Li Auto (理想) también existía la preocupación de que «usar modelos de código abierto podría perjudicar al equipo de desarrollo propio». La respuesta de Li Xiang es que un modelo de lenguaje de código abierto puede servir como base, pero lo que Li Auto quiere es desarrollar, a partir de sus escenarios de negocio, su propio VLA, su Agent OS y sus capacidades en el mundo físico. Lo central aquí no es elegir entre «código abierto o desarrollo propio», sino cómo se integra la base de código abierto en los datos, las tareas, la cadena de herramientas y los escenarios propios de la empresa.

\begin{knowledgebox}{Relación entre los modelos de código abierto y las capacidades de la empresa}
Linux puede ser de código abierto, y aun así Android formó sus propios componentes, API, ecosistema y experiencia de producto. Del mismo modo, los modelos de lenguaje pueden ser de código abierto, pero los fabricantes de automóviles siguen necesitando construir sus propias capacidades de sistema en torno al automóvil, los robots, los datos y los escenarios de los usuarios.
\end{knowledgebox}

\subsection{Resumen de la sección}

El capítulo inicial marca el tono de todo el episodio: la IA no es solo un contexto amplio y muchos parámetros, sino la manera de convertir la capacidad de procesar información en capacidad de actuar. DeepSeek se discute porque comprimió la investigación, la ingeniería y la organización en una metodología que se puede aprender.
\section{Experto técnico 1: VLA es el Driver OS del mundo físico}

La sección anterior trató los productos de IA y los métodos de organización; esta sección entra en el canal de expertos técnicos. Zhang Xiaojun (张小珺) invoca al «experto técnico del gran modelo CEO» y le pide a Li Xiang (李想) que hable de VLA. VLA significa Vision-Language-Action: Vision se encarga de percibir el mundo físico, Language de comprender reglas, intenciones y conocimiento, y Action de generar acciones. Para Li Auto (理想), VLA no es un juguete multimodal abstracto, sino un «gran modelo conductor», y también el gran modelo o sistema operativo más importante para el ámbito del automóvil y del transporte.

\begin{figure}[H]
\centering
\includegraphics[width=0.86\textwidth]{vla-driver-os-stack.png}
\caption{Sistema operativo conductor VLA: Vision, Language y Action se combinan en un gran modelo de conducción del mundo físico. Diagrama conceptual propio, elaborado a partir de la conversación de 00:36:18--00:42:30.}
\end{figure}

\begin{knowledgebox}{Lectura de la figura: las tres capas de VLA no son sustantivos paralelos}
Vision resuelve «ver», Language resuelve «comprender» y Action resuelve «hacer». Lo verdaderamente difícil en la conducción autónoma es integrar las tres en un ciclo cerrado que se pueda entrenar, evaluar y desplegar, no ensamblar por separado un modelo de visión, un modelo de lenguaje y un controlador.
\end{knowledgebox}

\subsection{De los algoritmos de reglas al enfoque de extremo a extremo, y de ahí a VLA}

Li Xiang describe la evolución de la conducción inteligente como un proceso evolutivo, no como una ruptura. Al principio hubo algoritmos basados en reglas; después, el enfoque de extremo a extremo más VLM; ahora se entra en la etapa de VLA. El cambio fundamental es que el sistema ya no solo «entiende» y «planifica», sino que debe incorporar la generación de acciones al aprendizaje. El objetivo de VLA no es que el auto aprenda a razonar en palabras, sino que en carreteras reales ejecute maniobras de conducción más parecidas a las de una persona, más estables y más seguras.

\begin{importantbox}{Definición técnica de VLA}
VLA es un modelo que integra la entrada visual, la representación de lenguaje y conocimiento, y la salida de acciones en una misma cadena de entrenamiento e inferencia. En la conducción autónoma, su salida no es texto, sino el comportamiento del vehículo; por eso su evaluación tampoco puede limitarse a la calidad de las respuestas, sino que debe considerar la seguridad, la comodidad, el cumplimiento de las reglas y el logro del objetivo.
\end{importantbox}

\subsection{Modelo base de 32B en la nube y despliegue en el vehículo}

En la entrevista se menciona que Li Auto está entrenando un modelo base VL de 32B en la nube. El modelo en la nube se encarga de procesar más datos, con más parámetros y un entrenamiento más complejo; el modelo en el vehículo, en cambio, está limitado por la capacidad de cómputo, la latencia y los requisitos de seguridad, y necesita pasar por destilación, despliegue y validación para llegar al vehículo. Esta estructura se corresponde con la fábrica de modelos en la nube de la que habló XPeng (小鹏) en el EP120: la IA del mundo físico no suele ser «un modelo que se sube al auto», sino una ingeniería de sistemas de «entrenamiento en la nube, compresión en el dispositivo y validación en condiciones reales».

\begin{knowledgebox}{Términos clave: cadena de entrenamiento de VLA}
\begin{tabularx}{\textwidth}{p{0.22\textwidth}p{0.32\textwidth}X}
\toprule
Etapa & Función & Riesgos y restricciones \\
\midrule
Modelo base VL & Hacer que el modelo comprenda visión y lenguaje & Requiere datos multimodales a gran escala y capacidad de cómputo. \\
Entrenamiento de Action & Convertir la comprensión en acciones de control & La salida afecta directamente la seguridad y la experiencia. \\
Datos de modelos del mundo & Proporcionar escenarios controlables de entrenamiento y prueba & Sigue habiendo una brecha entre la simulación y el mundo real. \\
Modelo en el dispositivo & Ejecutarse en tiempo real en el vehículo & Restricciones más estrictas de cómputo, latencia, consumo de energía y regulación. \\
Validación en carretera & Determinar la capacidad real & La validación es costosa y es difícil cubrir los escenarios de cola larga. \\
\bottomrule
\end{tabularx}
\end{knowledgebox}

\subsection{Resumen de la sección}

En este episodio, VLA no es una palabra de moda, sino un sistema de producto: debe integrar de principio a fin el modelo base en la nube, el modelo del mundo, el entrenamiento de acciones, el despliegue en el dispositivo y la validación en carretera. Solo así la IA puede pasar de «responder preguntas» a «conducir vehículos».
\section{Experto técnico II: modelo del mundo, retroalimentación por refuerzo y ciclo cerrado del tráfico}

La sección anterior explicó qué es VLA; esta sección explica cómo se entrena y se valida. Li Xiang (李想) mencionó que Li Auto (理想) usa un modelo del mundo para generar datos de entrenamiento: el modelo conduce del punto A al punto B, y el resultado se evalúa con retroalimentación sobre comodidad, colisiones, reglas de tránsito y otros aspectos. Aquí se aprecia un patrón muy claro: la flota real aporta los datos, el modelo del mundo genera escenarios y aplica los exámenes, VLA se actualiza con esos datos y esa retroalimentación, y al final se despliega de nuevo en los vehículos.

\begin{figure}[H]
\centering
\includegraphics[width=0.96\textwidth]{world-model-loop.png}
\caption{Ciclo cerrado del modelo del mundo del tráfico: los datos de conducción reales, la generación por simulación y el entrenamiento de VLA se impulsan mutuamente. Diagrama conceptual propio, elaborado a partir de la conversación en 00:47:30--00:48:10 y 01:10:06--01:10:48.}
\end{figure}

\begin{knowledgebox}{Lectura de la figura: el modelo del mundo cumple tres funciones}
Primero, puede servir como sistema de examen para evaluar el desempeño de VLA en distintos escenarios de tráfico. Segundo, puede generar datos de entrenamiento que cubren escenarios difíciles de recopilar con frecuencia en las vías reales. Tercero, en el futuro podría formar parte del sistema de operación de vehículos autónomos y ayudar a los vehículos sin conductor a entender y predecir el mundo del tráfico.
\end{knowledgebox}

\subsection{El modelo del mundo no es un «generador de imágenes»}

En el contexto de la conducción autónoma, el modelo del mundo no se limita a generar videos atractivos. Debe ser sensible a las consecuencias de las acciones: cómo evoluciona el mundo del tráfico si el vehículo cambia de carril, acelera, esquiva o espera. También debe ser sensible a la seguridad y a las reglas: las colisiones, los frenados bruscos, pasarse un semáforo en rojo, pisar las líneas del carril, colarse en la fila y la comodidad no son solo cuestiones visuales, sino cuestiones del resultado de la tarea.

\begin{importantbox}{La fórmula mínima del modelo del mundo}
El modelo del mundo del tráfico puede escribirse, de forma aproximada, como:
\[
s_{t+1} = f(s_t, a_t, c_t).
\]
Aquí, \(s_t\) representa el estado actual del tráfico, \(a_t\) la acción del vehículo, \(c_t\) las restricciones y reglas del escenario, y \(s_{t+1}\) el estado en el instante siguiente. Lo verdaderamente difícil es lograr que \(f\) sea lo bastante confiable en los escenarios de cola larga de las vías reales.
\end{importantbox}

\subsection{Costo de validación y escala del modelo}

Li Xiang mencionó que el costo de validación bajó de 180,000 yuanes a 4,000 yuanes por cada 10,000 kilómetros. Esto indica que la ruta de VLA y el modelo del mundo no es solo una cuestión de capacidad del modelo, sino que también cambia la economía de la validación. Si la conducción autónoma solo pudiera validarse acumulando kilometraje en vías reales, el costo sería altísimo; el modelo del mundo y la simulación pueden reducir ciertos costos de prueba, pero al final la validación en el mundo real sigue siendo indispensable como respaldo.

\begin{warningbox}{La simulación no puede sustituir directamente al mundo real}
El modelo del mundo puede reducir los costos de exploración y validación, pero no puede reemplazar por completo al mundo real. Los comportamientos de cola larga que la simulación no cubre, el ruido de los sensores, las obras viales, el clima extremo y las acciones irracionales de las personas pueden seguir apareciendo en las vías reales.
\end{warningbox}

\subsection{L3, L4 y el límite de cómputo a bordo del vehículo}

Ya se explicó cómo el modelo del mundo reduce los costos de entrenamiento y validación; esta subsección lleva el problema a los niveles de conducción autónoma y a las restricciones del cómputo a bordo. En la entrevista, Li Xiang relacionó la capacidad L3/L4 con la escala de los modelos en la nube y a bordo. A su juicio, la generación actual de cómputo corresponde aproximadamente a la capacidad L3, y L4 todavía requiere más capacidades y más condiciones. El valor didáctico de esta apreciación es que el nivel de conducción autónoma no es un simple interruptor de software, sino que lo determinan en conjunto la escala del modelo, el cómputo a bordo, el sistema de validación, la regulación y la responsabilidad operativa.

\begin{knowledgebox}{Términos clave: L3 y L4}
L3 suele significar que, en condiciones específicas, el sistema puede asumir la tarea de conducción, pero una persona debe tomar el control cuando el sistema lo solicita; L4 significa que, dentro de un dominio de diseño operativo delimitado, el sistema puede funcionar sin depender de que una persona tome el control. La diferencia entre ambos no está solo en las puntuaciones del modelo, sino en los límites de responsabilidad, el diseño de redundancia, la regulación y el sistema de operación.
\end{knowledgebox}

\subsection{Resumen de la sección}

El modelo del mundo lleva a VLA de ser «un modelo que sabe conducir» a ser «un sistema que puede entrenarse, examinarse y operarse». Reduce el costo de validación, pero no elimina la incertidumbre del mundo real; eleva el techo de desempeño, pero también exige más cómputo, más datos y un sistema de seguridad más sólido.
\section{De responder a actuar: Agent OS y exigencias de alineación}

Esta sección extiende el problema de la conducción autónoma a los Agents. Li Xiang (李想) considera que en cinco años no necesariamente habrá un Agent general, pero sí habrá un Agent OS. Este juicio significa lo siguiente: cada oficio profesional necesita un Agent distinto, pero todos esos Agents deben operar sobre una base común similar a un sistema operativo, de la cual obtienen herramientas, permisos, contexto, máquinas virtuales, memoria y límites de seguridad. Un Agent OS no es una ventana de chat, sino la infraestructura de los Agents profesionales.

\begin{figure}[H]
\centering
\includegraphics[width=0.96\textwidth]{action-agent-alignment.png}
\caption{De responder a actuar: cuando la IA interviene sobre activos, herramientas y vehículos, las exigencias de alineación aumentan. Diagrama conceptual propio, elaborado a partir de la conversación de 00:53:58--00:55:00 y 00:55:00--00:55:25.}
\end{figure}

\begin{knowledgebox}{Lectura de la figura: la acción amplifica el riesgo}
Si una búsqueda de información falla, se puede buscar de nuevo; si una investigación tiene errores, se puede revisar. Pero si un Agent opera activos, herramientas, computadoras, vehículos o bienes, el error modifica el mundo de forma directa. Por eso las exigencias de alineación, los límites de permisos y la atribución de responsabilidades de un Agent son mayores que las de un modelo conversacional.
\end{knowledgebox}

\subsection{Por qué un Agent general rinde menos que un Agent profesional}

Li Xiang sostiene que las herramientas de producción reales deben sustituir y liberar a las personas de su trabajo de alta frecuencia, como conducir, programar, operar, vender o fabricar. Cada oficio tiene sus propios datos, herramientas, permisos, costos de error y criterios de evaluación; por ello, a corto plazo es más probable que surjan Agents profesionales que un Agent capaz de todo. El valor del Agent OS consiste en ofrecer a esos Agents profesionales un entorno de ejecución común.

\begin{importantbox}{Definición de Agent OS}
Un Agent OS es el sistema base sobre el que operan los Agents profesionales. Incluye, como mínimo, gestión de contexto, invocación de herramientas, control de permisos, máquinas virtuales o entornos de ejecución, memoria, registros, evaluación y políticas de seguridad. Su valor no está en «poder hacer de todo», sino en permitir que distintos Agents profesionales hagan su trabajo de forma confiable.
\end{importantbox}

\subsection{La alineación pasa de ser un problema de lenguaje a un problema de acción}

Cuando la IA solo responde preguntas, la alineación se refiere sobre todo a la seguridad del contenido, la exactitud factual y las preferencias de valores; cuando la IA empieza a invocar herramientas y activos, la alineación se convierte en riesgo operativo. Por ejemplo, el error de un Agent de conducción afecta la seguridad vial; el de un Agent financiero afecta los fondos; el de un Agent de código afecta la estabilidad del sistema. Así, la alineación se amplía de «si el modelo dice lo correcto» a «si lo que el modelo hace es controlable».

\begin{knowledgebox}{Estratificación aproximada de las exigencias de alineación}
\[
\text{risk} \approx \text{action scope} \times \text{asset value} \times \text{irreversibility}.
\]
Aquí, \(\text{action scope}\) indica el alcance de lo que el Agent puede operar, \(\text{asset value}\) indica el valor de los activos invocados e \(\text{irreversibility}\) indica qué tan difícil es deshacer un error. Mientras más cerca se esté de vehículos, dinero, cuentas, sistemas de producción y dispositivos físicos, mayores serán las exigencias de alineación.
\end{knowledgebox}

\subsection{Resumen de la sección}

El Agent OS es el puente que, en este episodio, lleva de la conducción autónoma a los productos de IA de uso general. Muestra que el siguiente paso de la IA no es conversar mejor, sino actuar mejor en entornos con permisos, herramientas y responsabilidades.
\section{Experto en estrategia I: escala, necesidades del usuario, tecnología y producto, capacidad organizacional}

El capítulo anterior trató el sistema de acción; este capítulo entra en el canal del experto en estrategia. Al hablar de la reunión estratégica de 2025 en Yanqi Hu (雁栖湖), Li Xiang (李想) sitúa la estrategia en un marco de cuatro variables: en el centro está la escala y, alrededor, las necesidades del usuario, la tecnología y el producto, y la capacidad organizacional. La escala no es un simple número de ingresos, sino algo que modifica la estructura de usuarios, la complejidad del producto, la forma de colaborar dentro de la organización y el panorama competitivo externo. Cuando una empresa pasa de ingresos de diez mil millones a cien mil millones, y después avanza hacia una escala aún mayor, tiene que volver a aprender las mejores prácticas del sector.

\begin{figure}[H]
\centering
\includegraphics[width=0.96\textwidth]{org-as-moe.png}
\caption{La organización también funciona como un MoE: el CEO asigna expertos, metodologías, talento joven y ciclos de retroalimentación. Diagrama conceptual propio, elaborado a partir de la conversación de 01:25:36--02:03:20.}
\end{figure}

\begin{knowledgebox}{Lectura de la figura: la organización no es una suma de personas, sino la asignación de expertos}
La tecnología, el producto, la estrategia, la contratación de recién egresados, la metodología y la retroalimentación forman en conjunto la capacidad organizacional. La organización se parece a un MoE porque las capacidades de los distintos expertos deben asignarse correctamente; si el enrutamiento es erróneo, por muchos expertos que haya, sus esfuerzos se anulan entre sí.
\end{knowledgebox}

\subsection{De Toyota, GM, Google y Huawei a una metodología propia}

El apartado anterior presentó la organización como un sistema de asignación de expertos; esta sección vuelve a cómo se forma ese sistema. En sus primeros años, Li Auto (理想) estudió el sistema de trabajo de Toyota, el proceso de investigación y desarrollo de GM y los OKR de Google; más adelante analizó también a Huawei y a Apple. Este aprendizaje no consistió en copiar, sino en buscar métodos transferibles para cada etapa de escala. En la etapa del Li ONE (理想 ONE) lo necesario era entregar; en la etapa de ingresos de cien mil millones se requieren densidad de talento, coordinación organizacional y capacidades de plataforma; en la etapa de dispositivos terminales de la era de la AGI se necesitan modelos, sistema operativo, hardware y ecosistema.

\begin{warningbox}{La metodología no debe tratarse como una religión}
Aprender de Toyota, Google, Huawei o Apple no significa copiar el lema organizacional de alguna de ellas. La metodología debe servir a la escala actual, a la etapa del producto y a las carencias de capacidad. Cuando cambia la escala de la empresa, también cambian las necesidades del usuario y la forma de la organización, y un método antiguo puede pasar de ser un activo a ser una carga.
\end{warningbox}

\subsection{Organizaciones pequeñas y alta densidad de talento}

Si la metodología determina la forma de colaborar, la siguiente pregunta es quién la ejecuta. Li Xiang menciona en varias ocasiones los equipos pequeños, la alta densidad de talento y la proporción de contrataciones de recién egresados. Un cambio posible en la era de la IA es que más organizaciones pequeñas obtengan mayor capacidad intelectual y fuerza anímica mediante modelos, herramientas y metodologías, en lugar de crecer linealmente en número de personas. El ejemplo de DeepSeek también se usa para ilustrarlo: el equipo no tiene que ser grande, pero sí debe tener alta densidad de talento, una dirección unificada y un ciclo de ingeniería sólido de principio a fin.

\begin{importantbox}{Los cuatro niveles de la capacidad organizacional}
El primer nivel es la densidad de talento: si las personas son lo bastante capaces. El segundo es la metodología: si las personas capaces colaboran con un mismo lenguaje. El tercero es la retroalimentación: si el negocio real y los usuarios corrigen continuamente el juicio. El cuarto es la energía: si el equipo está dispuesto a sostener tareas difíciles de manera continua.
\end{importantbox}

\subsection{Resumen de la sección}

El canal del experto en estrategia entiende a Li Auto dentro de los cambios de escala. Una empresa no avanza linealmente del punto A al punto B, sino que, cada vez que supera un escalón de escala, debe reconfigurar las necesidades del usuario, la tecnología y el producto, y la capacidad organizacional.
\section{Segundo experto en estrategia: el terminal de la era de la AGI y la próxima Apple}

Esta sección trata la visión estratégica más ambiciosa de la entrevista: si Li Auto (理想) puede convertirse en la Apple de la era de la AGI. Li Xiang (李想) define el terminal de la era de la AGI como aquel que reúne cuatro tipos de capacidades: percepción del mundo físico en 360 grados, cognición y toma de decisiones, Action, y reflexión y retroalimentación. Esta definición amplía el «terminal» de los dispositivos con pantalla a los puntos de acceso al mundo físico: el automóvil puede ser un terminal, el robot también, y los dispositivos domésticos y los equipos industriales también podrían llegar a serlo.

\begin{figure}[H]
\centering
\includegraphics[width=0.96\textwidth]{agi-era-apple-route.png}
\caption{La ruta hacia la Apple de la era de la AGI: el automóvil, los robots, el sistema operativo y las capacidades del modelo forman en conjunto el punto de acceso. Diagrama conceptual propio, elaborado a partir de la conversación de 01:38:20--01:46:40 y 02:09:16--02:09:21.}
\end{figure}

\begin{knowledgebox}{Lectura de la figura: el automóvil no es el único destino, sino el punto de acceso físico}
El automóvil, el VLA, los robots, el sistema operativo y AGI Apple de la figura no forman un lema lineal. El automóvil es el punto de acceso físico de mayor escala que existe hoy, el VLA es el cerebro de la conducción, los robots extienden el alcance al hogar y a la fábrica, y el sistema operativo conecta el software con las herramientas; solo al final podría formarse la plataforma de la siguiente generación.
\end{knowledgebox}

\subsection{Capacidades de software: sistema operativo, máquinas virtuales y herramientas}

La sección anterior situó al automóvil y a los robots dentro de la visión de la plataforma de la siguiente generación; esta sección desglosa primero la base de software. Li Xiang considera que el terminal de la era de la AGI requiere distintas capacidades de software. Primero, capacidad de sistema operativo, que permita a la IA funcionar de manera estable en el mundo físico y en el digital. Segundo, capacidad de máquina virtual o de entorno de ejecución, para que el Agent funcione entre la computadora local, la nube y los dispositivos. Tercero, un ecosistema de herramientas, porque el Agent necesita invocar herramientas y las propias herramientas también se rediseñarán para la IA.

\begin{importantbox}{Por qué el sistema operativo vuelve a ser importante}
Cuando la IA es solo una aplicación, el sistema operativo es como un telón de fondo; cuando la IA debe observar de forma continua, invocar herramientas, controlar dispositivos, registrar estados y ejecutar acciones, el sistema operativo se convierte en el núcleo de la seguridad, el rendimiento, los permisos y el ecosistema.
\end{importantbox}

\subsection{Capacidades de hardware: percepción, capacidad de cómputo y cuerpo físico}

La base de software resuelve «cómo se ejecuta»; esta sección pasa a «en qué cuerpo se ejecuta». El terminal de la AGI también necesita capacidades de hardware. El automóvil tiene sensores, capacidad de cómputo, energía, espacio y movilidad; el robot tiene cuerpo físico, actuadores, tacto y escenarios de fábrica u hogar; los dispositivos domésticos podrían necesitar una percepción unificada y un cerebro unificado. En la entrevista, Li Xiang no presenta la ruta de los robots como la única respuesta, sino que reconoce dos posibilidades: los robots humanoides pueden usar herramientas humanas, o bien se pueden transformar las herramientas y el entorno para que dispositivos no humanoides completen las tareas.

\begin{knowledgebox}{Dos visiones de la ruta de los robots}
\begin{tabularx}{\textwidth}{p{0.25\textwidth}p{0.34\textwidth}X}
\toprule
Ruta & Ventaja & Dificultad \\
\midrule
Robot humanoide & Se adapta a las herramientas y espacios humanos existentes & Cuerpo físico complejo, costo alto, control difícil. \\
Transformar el entorno o las herramientas & Es más fácil lograr alta confiabilidad en una sola tarea & Exige rediseñar el ecosistema del hogar, de la fábrica o de los dispositivos. \\
Cerebro unificado & Permite reutilizar el modelo, la memoria y la comprensión de tareas & Los espacios de acción difieren mucho entre cuerpos físicos distintos. \\
\bottomrule
\end{tabularx}
\end{knowledgebox}

\subsection{El desfase entre las ventas y el valor de la IA}

El mercado todavía evalúa a Li Auto principalmente por sus ventas, porque el valor de servicios en la nube, el valor de software y el valor de plataforma que aportan las capacidades de IA aún no se reflejan por completo. Este desfase no es raro: en las primeras etapas de una transición de plataforma, los observadores externos suelen evaluar las capacidades nuevas con indicadores viejos. El verdadero punto de inflexión llegará cuando las capacidades de IA mejoren la experiencia, reduzcan costos, generen ingresos o den lugar a nuevas formas de terminal.

\begin{warningbox}{«Convertirse en Apple» no es una metáfora de marca, sino de plataforma}
Si se entiende a Apple solo como una marca de gama alta, esta discusión se malinterpreta. Li Xiang habla de capacidades de plataforma: hardware, sistema operativo, ecosistema de desarrollo, punto de acceso de usuarios, cadena de herramientas y modelo de negocio. Para que un fabricante de automóviles se convierta en la Apple de la era de la AGI, debe demostrar que no se limita a vender automóviles.
\end{warningbox}

\subsection{Resumen de la sección}

La visión del terminal de la era de la AGI amplía la estrategia de Li Auto del automóvil a la plataforma. El automóvil es el punto de acceso, el VLA es la capacidad, el Agent OS es la capa de software y los robots son la extensión; la cuestión de fondo es si la empresa puede convertir el punto de acceso al mundo físico en una plataforma inteligente.
\section{Organización y personas: ir contra la naturaleza humana, seguirla y el cerebro común}

La sección anterior trató la plataforma; esta sección vuelve a la organización. Li Xiang (李想) mencionó en repetidas ocasiones la idea de «hablar siguiendo la naturaleza humana, actuar yendo contra ella». La frase puede entenderse fácilmente como una técnica de gestión, pero en el contexto de la entrevista apunta a la dificultad del cambio organizacional: las mejores prácticas, el desarrollo de capacidades, la visión de largo plazo, la revisión retrospectiva y la autocrítica suelen ir contra la naturaleza humana; en cambio, hacer lo que uno quiere, evadir en el corto plazo y justificarse a sí mismo la siguen con más facilidad. Si una organización quiere evolucionar de forma continua, debe diseñar un mecanismo que comprenda la naturaleza humana y que, al mismo tiempo, impulse un crecimiento que va en su contra.

\begin{knowledgebox}{Ir contra la naturaleza humana no es reprimir a las personas, sino combatir la inercia}
Cuando Li Xiang habla de «ir contra la naturaleza humana», se refiere a combatir la pereza, la evasión y la comodidad de corto plazo, no a negar los sentimientos de las personas. Una buena organización debe comunicarse respetando la forma de expresarse y la dignidad de las personas, pero no puede ser permisiva en el desarrollo de capacidades ni en la exigencia de resultados.
\end{knowledgebox}

\subsection{El cerebro común y el corazón común de tres a siete personas}

El apartado anterior explicó que la organización debe comprender la naturaleza humana sin consentir la inercia; este apartado examina la unidad organizacional concreta que propone Li Xiang. En la segunda mitad de la entrevista, habló de un modelo: de tres a siete personas forman un cerebro más fuerte, un corazón más fuerte y una energía más fuerte. No se trata de una discusión de grupo ordinaria, sino de reunir la complementariedad de capacidades, la energía emocional, la responsabilidad compartida y una metodología unificada. Si una persona es muy capaz pero no obtiene resultados, puede deberse a que le falta esta estructura común.

\begin{importantbox}{Cuatro condiciones del cerebro común}
Primera, los miembros se complementan entre sí en lugar de duplicarse. Segunda, comparten un objetivo y una metodología. Tercera, la relación es lo bastante cercana para dar retroalimentación sincera. Cuarta, el equipo tiene la energía para soportar la presión en conjunto, y no solo una tabla de reparto de tareas.
\end{importantbox}

\subsection{Preocuparse por los usuarios y por las personas cercanas}

Li Xiang resume la esencia de la conexión entre personas en dos formas de «preocuparse»: preocuparse por los usuarios y preocuparse por las personas cercanas. Preocuparse por los usuarios es un consenso de valores; preocuparse por las personas cercanas es la base de la colaboración. Si una organización solo atiende los asuntos y no a las personas, desgasta fácilmente las relaciones; si solo atiende las relaciones y no a los usuarios, pierde los resultados. Una buena organización debe mantener la tensión entre ambos aspectos.

\begin{warningbox}{«Juzgar los hechos y no a las personas» a veces encubre los problemas reales}
Las organizaciones de ingeniería suelen decir que se juzgan los hechos y no a las personas, lo cual evita ataques emocionales. Pero si se deja de mirar por completo a las personas, se pasan por alto su crecimiento, su motivación, su energía y la calidad de sus relaciones. Cuando Li Xiang subraya aquí «primero las personas, luego los asuntos», recuerda que los problemas de una organización se reducen, en última instancia, a si las personas se vuelven más fuertes.
\end{warningbox}

\subsection{Resumen de la sección}

La capacidad organizacional no es un diagrama de flujo, sino la manera en que las personas crecen juntas bajo una presión prolongada. En la era de la IA, las organizaciones pequeñas pueden ser más fuertes, siempre que cuenten con un cerebro común, un corazón común, una metodología común y una energía común.
\section{Energía, relaciones íntimas y programa de memoria}

El capítulo anterior entendía la organización como un cerebro común y un corazón común; este capítulo pasa a la fuente de energía, en un nivel más profundo. Esta sección examina primero la parte de la entrevista que menos parece técnica, pero que mejor explica la metodología de Li Xiang (李想): la energía proviene del crecimiento, de las relaciones íntimas, de ser necesitado y de necesitar a otros. Dice que necesita a sus hijos, a su pareja y a los responsables de los departamentos de primer nivel, e incluso que ellos son más importantes para él que él para ellos. Esta afirmación no es sentimentalismo motivacional, sino una forma de ver las relaciones como un sistema de retroalimentación: las relaciones íntimas permiten ver las propias carencias y también ayudan a mejorar de forma continua.

\begin{figure}[H]
\centering
\includegraphics[width=0.96\textwidth]{energy-relationship-loop.png}
\caption{Energía y relaciones íntimas: las relaciones no son un gasto, sino un sistema de retroalimentación que hace mejores a las personas. Diagrama conceptual propio, elaborado a partir de la conversación de 02:09:26--02:28:20.}
\end{figure}

\begin{knowledgebox}{Lectura de la figura: la energía no es entusiasmo, sino retroalimentación continua}
Necesidad, conexión, retroalimentación, energía y crecimiento forman un ciclo. Las relaciones íntimas importan porque ofrecen retroalimentación real; la retroalimentación real expone las carencias; solo al aceptar las carencias y mejorar se genera energía sostenida.
\end{knowledgebox}

\subsection{Los límites de las relaciones íntimas}

Antes se subrayó que las relaciones íntimas aportan energía; esta sección añade una condición de límite igual de importante. Li Xiang también insiste en no construir demasiadas relaciones íntimas. Una relación íntima no es la vida social en general, sino los familiares directos, unos pocos amigos de muchos años y las personas con quienes se comparte la responsabilidad en el trabajo. Lo que puede causar daño suelen ser también solo las relaciones íntimas, porque son las que de verdad llegan a las necesidades y a los puntos vulnerables de una persona. Esta idea puede trasladarse a la organización: las relaciones de alta calidad aportan energía; las de baja calidad la consumen.

\begin{importantbox}{Definición de relación íntima}
Una relación íntima no consiste en conversar con frecuencia ni en tener cercanía social, sino en que ambas partes puedan necesitarse de verdad, darse retroalimentación real, asumir juntas las consecuencias y, a largo plazo, hacer mejor a la otra.
\end{importantbox}

\subsection{Programa de memoria: si pudieras modificar un fragmento del código de tu vida}

En la entrevista también surgió la pregunta del «programa de memoria»: si se pudieran borrar o modificar recuerdos, qué se elegiría. La respuesta de Li Xiang no considera simplemente malos los recuerdos dolorosos, porque muchos recuerdos moldearon la metodología, las relaciones y el juicio que tiene hoy. Aquí la «memoria» puede entenderse como los datos de entrenamiento de una persona: no es solo un archivo, sino también la fuente de sus estrategias de conducta y de sus juicios de valor.

\begin{knowledgebox}{La memoria como datos de entrenamiento de las personas}
Un modelo de IA forma sus parámetros a partir de los datos; una persona forma su metodología a partir de sus recuerdos, sus relaciones y sus acciones. Borrar los recuerdos dolorosos parece reducir el desgaste, pero también podría borrar una fuente de capacidad. Lo importante es asimilar los recuerdos hasta convertirlos en método, en lugar de que los recuerdos se activen una y otra vez.
\end{knowledgebox}

\subsection{Resumen de la sección}

La energía y las relaciones íntimas no son un «interludio humanista» en la entrevista, sino variables de fondo con las que Li Xiang entiende la organización y la inteligencia. Sin energía, la metodología no puede sostenerse a largo plazo; sin relaciones reales, es difícil que una persona vea sus puntos ciegos.
\section{La sabiduría es nuestra relación con todas las cosas}

Esta sección culmina en la definición central de la segunda mitad de la entrevista: la sabiduría es nuestra relación con todas las cosas. Esta definición es singular porque no entiende la sabiduría como IQ, cantidad de conocimiento o puntaje de razonamiento, sino como la calidad de las relaciones entre personas, entre las personas y las cosas, las organizaciones, la tecnología y el mundo. Por ello, el sentido de la IA tampoco se limita a sustituir a las personas: consiste en liberarlas de acciones de bajo valor que consumen energía, para que dispongan de más energía con la cual atender relaciones de mayor calidad.

\begin{figure}[H]
\centering
\includegraphics[width=0.92\textwidth]{wisdom-relation-map.png}
\caption{La sabiduría es relación: la sabiduría surge del contacto continuo con las personas, las cosas, las organizaciones, la tecnología y el mundo. Diagrama conceptual propio, elaborado a partir de la conversación de 02:26:47--02:31:10.}
\end{figure}

\begin{knowledgebox}{Lectura de la figura: la sabiduría no es un puntaje mental}
Las personas, las cosas, las organizaciones, la tecnología y el mundo constituyen en conjunto la fuente de la sabiduría. Si una persona solo razona dentro de su cabeza, sin entrar en contacto con personas y hechos reales, su sabiduría se empobrece; si una IA solo circula entre tokens, sin pasar a la acción y a la retroalimentación, también le resulta difícil generar valor real.
\end{knowledgebox}

\subsection{El sentido de la IA: reducir las acciones que consumen energía}

El apartado anterior definió la sabiduría como relación; este apartado se ocupa de cómo la IA cambia la distribución de las acciones dentro de esas relaciones. Li Xiang (李想) pone el ejemplo del personal de ventas: a los buenos vendedores y especialistas de producto les gusta tratar con los clientes, porque ese contacto les da energía; en cambio, las invitaciones repetitivas, los procesos de bajo valor y la comunicación mecánica pueden agotarla. Uno de los valores de la IA es automatizar estas acciones que consumen energía, para que las personas dediquen su energía a tareas con más sabiduría y mayor calidad de relación.

\begin{importantbox}{La relación entre la IA y las personas}
La IA no solo sustituye trabajo: también redistribuye la energía de las personas. Un sistema de IA verdaderamente bueno debería reducir las acciones de bajo valor, repetitivas y desgastantes, y al mismo tiempo conservar o fortalecer el contacto de alta calidad de las personas con los clientes, los productos, las organizaciones y el mundo.
\end{importantbox}

\subsection{Del envoltorio a las capacidades básicas}

Si el sentido de la IA está en asumir más action, al final es inevitable volver a la cuestión de las capacidades básicas. Hacia el cierre, Li Xiang expone su opinión sobre aplicaciones Agent como «Xiaohong» (小红). Considera que llevar la IA hacia la action real, la action profesional y la to-do list es un intento importante; pero si los competidores son empresas con modelos sólidos, no basta con un envoltorio, máquinas virtuales y capacidad de uso de herramientas: las capacidades necesarias del modelo hay que entrenarlas de verdad. Este planteamiento deja una advertencia muy clara para los emprendimientos de aplicaciones de IA: el punto de entrada, los flujos y la experiencia de producto pueden ponerse en marcha primero, pero al final hay que completar las capacidades básicas.

\begin{figure}[H]
\centering
\includegraphics[width=0.94\textwidth]{shell-vs-model-ability.png}
\caption{Envoltorio frente a capacidad del modelo: la innovación en aplicaciones debe avanzar hacia la action y, en última instancia, completar las capacidades básicas. Diagrama conceptual propio, elaborado a partir de la conversación de 02:44:21--02:45:30.}
\end{figure}

\begin{knowledgebox}{Lectura de la figura: el envoltorio tiene valor, pero no hay que quedarse ahí}
A la izquierda, el envoltorio y las herramientas permiten validar con rapidez las necesidades de los usuarios y los flujos de trabajo; a la derecha, las capacidades básicas determinan el techo a largo plazo. Si una aplicación Agent quiere llegar a la action profesional, tiene que completar de manera gradual el modelo, la memoria, la invocación de herramientas, la evaluación y la confiabilidad.
\end{knowledgebox}

\begin{warningbox}{No interpretar «el envoltorio no es innovación» como una postura contra las aplicaciones}
La innovación en aplicaciones es muy importante, porque es la más cercana a los usuarios y a los escenarios de acción. Sin embargo, cuando el escenario pasa a tareas de alto valor, con fuerte exigencia de alineación y alta confiabilidad, depender solo de modelos externos y de flujos de trabajo delgados choca con un techo, y las capacidades básicas tienen que estar a la altura.
\end{warningbox}

\subsection{Resumen de la sección}

«La sabiduría es relación» devuelve todo el episodio de la tecnología a las personas. VLA, Agent OS, los modelos del mundo y la metodología organizacional no buscan, en última instancia, lucirse técnicamente, sino que la inteligencia se integre mejor en la relación entre las personas y el mundo.
\section{Síntesis y ampliación}

Esta sección condensa toda la entrevista en unas cuantas conclusiones reutilizables. Primero, los productos de IA pasan de ser herramientas de información a ser herramientas de producción, y lo decisivo es la action: si pueden sustituir y liberar trabajo real, y si pueden entrar en herramientas, activos, vehículos y equipos. Segundo, el núcleo de la IA del mundo físico no es un modelo único, sino un ciclo cerrado formado por VLA, modelos del mundo, cómputo en el dispositivo, sistemas de verificación y responsabilidad regulatoria. Tercero, en la era de la AGI el dispositivo final no es solo el teléfono o el automóvil, sino una plataforma capaz de percibir el mundo físico, tomar decisiones cognitivas, ejecutar acciones y reflexionar sobre la retroalimentación. Cuarto, en la era de la IA la capacidad organizacional resulta aún más importante, porque los modelos reducen ciertos costos de ejecución, pero aumentan la importancia del juicio, la metodología y la densidad de talento. Quinto, la sabiduría humana no es solo capacidad cognitiva, sino capacidad de relación.

\begin{importantbox}{EP118 dentro de la serie de entrevistas sobre IA de Zhang Xiaojun (张小珺)}
EP120 trata de cómo Xiaopeng (小鹏) reorienta la conducción autónoma hacia la Physical AI, EP121 trata de la robótica y los modelos del mundo de DeepMind, y EP119 hace una arqueología de la arquitectura de Attention; EP118, en cambio, lleva estas líneas técnicas a la perspectiva del CEO de una empresa: cómo el modelo se convierte en dispositivo final, cómo el dispositivo final se convierte en plataforma y cómo la plataforma, a su vez, exige que cambien la organización y las personas.
\end{importantbox}

\subsection{Takeaways técnicos}

\begin{enumerate}
\item Lo importante de VLA no es la etiqueta de multimodal, sino integrar la visión, el lenguaje y la acción en un ciclo cerrado unificado de entrenamiento y verificación.
\item El valor del World Model no está solo en generar datos, sino también en la evaluación, la verificación y la operación futura.
\item Un Agent OS es una vía de implementación a corto plazo más viable que un Agent general, porque el trabajo real requiere herramientas especializadas, permisos y límites de seguridad.
\item Cuando la IA entra en la action, el riesgo escala de errores de contenido a errores sobre activos, herramientas y el mundo físico.
\item La capa de aplicaciones puede avanzar primero, pero a largo plazo tiene que completar el modelo, la memoria, la evaluación y la confiabilidad de la ejecución.
\end{enumerate}

\subsection{Takeaways de gestión}

\begin{enumerate}
\item Las mejores prácticas suelen ir contra la naturaleza humana; la organización debe diseñar mecanismos que contrarresten la comodidad de corto plazo.
\item Un equipo pequeño puede ser muy fuerte, pero a condición de tener alta densidad de talento, objetivos compartidos, una metodología unificada y retroalimentación real.
\item Los cambios de escala modifican las necesidades de los usuarios, los productos técnicos y la capacidad organizacional; no se pueden resolver los problemas de una etapa nueva con los métodos de una etapa anterior.
\item Las relaciones íntimas y la energía no son temas menores: determinan si una persona y una organización pueden sostener tareas complejas a largo plazo.
\item La sabiduría puede entenderse como calidad de las relaciones: cuanto más reales son las relaciones entre personas, entre personas y objetos, entre personas y tecnología, y entre personas y el mundo, más sólido es el juicio.
\end{enumerate}

\subsection{Lecturas adicionales}

\begin{itemize}
\item Contrastar con la entrevista de EP120 a Liu Xianming (刘先明), de Xiaopeng, permite comparar cómo distintos fabricantes de automóviles plantean la Physical AI, VLA, la fábrica de modelos en la nube y el despliegue en el dispositivo.
\item Contrastar con la revisión de la arquitectura de Attention de EP119 permite ver cómo la arquitectura base de los modelos influye en el costo de inferencia del contexto largo, los Agent y la IA del mundo físico.
\item Contrastar con la entrevista de EP121 a Tan Jie (谭捷), de DeepMind, permite entender VLA y los modelos del mundo dentro de los problemas robóticos de transferencia entre cuerpos, datos de acción y despliegue real.
\end{itemize}

\end{document}
